Tables~\ref{tab:case-study-transformations} and
\ref{tab:case-study-resources} join the opening/memory case study through frozen
file identities. The full experiment retains 288 opening/saving attempts
across six variants and three repetitions. The selected
Enron workbook has a separate eighteen-attempt follow-up on the larger-memory
host, described below. Its measurements retain their own host and configuration.
The tables display Migration and QT/RR, with Enron size evidence and a separate
resource follow-up.
The payload table uses exact retained comparator outputs and separately regenerated
Sheet Shrink controls with verified native-candidate member identities. Per-trial
member inventories of timed fresh outputs were not retained. Native pixel counts
are reported only for measured, retained image pairs with matched origins.
Missing scopes and image-size differences remain explicit.

The compact resource table reports total elapsed time and library working-set
peaks. The artifact separately retains pipeline working sets and job-commit
peaks. Each displayed time is a median of completed attempts,
with memory maxima over all three attempts. Stopped attempts retain their
elapsed time and memory boundary. Existing prepared-output
variants measure subsequent library processing. Fresh Sheet Shrink rows include
cleanup within the measured pipeline. The common host is a Ryzen 5 4500U laptop
with six logical processors, 8 GiB installed RAM and Windows 11. Each job has a
4 GiB commit cap and a 900 s timeout, with the free-RAM reserve disabled.
The Sheet Shrink executable is frozen at revision \texttt{765c480} and uses six
Rayon and six writer workers. Native visual observations use the retained
paired-rebuild protocol.

\begin{table*}[t]
\centering
\footnotesize
\setlength{\tabcolsep}{2.5pt}
\renewcommand{\arraystretch}{1.08}
\caption{Selected-case uncompressed payload reduction and retained native pixel evidence. Payload is the sum of all ZIP-member sizes in decimal MB. Image-pair counts state measured/recorded coverage.}
\label{tab:case-study-transformations}
\begin{tabular}{lrl}
\hline
Variant & Payload MB (reduction) & Changed / compared pixels; coverage \\
\hline
\multicolumn{3}{l}{\textbf{Destatis Migration}} \\
Original & \CasePairBars[4.4]{248.39 (0\%)}{1.09990}{1.09990}{casebarblue} & Reference \\
SheetShrink\textsuperscript{*} & \CasePairBars[4.4]{16.46 (93\%)}{1.09990}{0.07288}{casebarblue} & \begin{tabular}[c]{@{}l@{}}0 / 101,874,379\\37/38 image pairs\\5 size differences\\Normalized baseline\end{tabular} \\
NXPowerLite & \CasePairBars[4.4]{197.48 (20\%)}{1.09990}{0.87447}{casebarblue} & Failed \\
XLStylesTool & \CasePairBars[4.4]{248.42 (-0.0095\%)}{1.09990}{1.10000}{casebarorange} & Failed \\
Check Performance & \CasePairBars[4.4]{20.15 (92\%)}{1.09990}{0.08922}{casebarblue} & Memory limit \\
Inquire & \CasePairBars[4.4]{116.44 (53\%)}{1.09990}{0.51559}{casebarblue} & Failed \\
\hline
\multicolumn{3}{l}{\textbf{Data.gov QT/RR}} \\
Original & \CasePairBars[4.4]{1350.77 (0\%)}{1.10000}{1.10000}{casebarblue} & Reference \\
SheetShrink\textsuperscript{*} & \CasePairBars[4.4]{29.91 (98\%)}{1.10000}{0.02436}{casebarblue} & No capture regions \\
NXPowerLite & \CasePairBars[4.4]{29.98 (98\%)}{1.10000}{0.02442}{casebarblue} & No capture regions \\
XLStylesTool & \CasePairBars[4.4]{1350.77 (-7.8e-05\%)}{1.10000}{1.10000}{casebarorange} & No capture regions \\
Check Performance & \CasePairBars[4.4]{25.59 (98\%)}{1.10000}{0.02084}{casebarblue} & No capture regions \\
Inquire & \CasePairBars[4.4]{30.13 (98\%)}{1.10000}{0.02454}{casebarblue} & No completed report \\
\hline
\multicolumn{3}{l}{\textbf{Enron b2d0dfef}} \\
Original & \CasePairBars[4.4]{353.15 (0\%)}{1.09999}{1.09999}{casebarblue} & Reference \\
SheetShrink\textsuperscript{*} & \CasePairBars[4.4]{2.58 (99\%)}{1.09999}{0.00804}{casebarblue} & No retained pairs (0/2824) \\
NXPowerLite & \CasePairBars[4.4]{2.55 (99\%)}{1.09999}{0.00795}{casebarblue} & \begin{tabular}[c]{@{}l@{}}6,437 / 11,640,701,446\\1033/1033 image pairs\end{tabular} \\
XLStylesTool & \CasePairBars[4.4]{353.16 (-0.0011\%)}{1.09999}{1.10000}{casebarorange} & No capture regions \\
Check Performance & \CasePairBars[4.4]{2.54 (99\%)}{1.09999}{0.00792}{casebarblue} & \begin{tabular}[c]{@{}l@{}}456,503 / 46,890,576\\8/1032 image pairs\\2 size differences\end{tabular} \\
Inquire & \CasePairBars[4.4]{2.57 (99\%)}{1.09999}{0.00801}{casebarblue} & \begin{tabular}[c]{@{}l@{}}9,739 / 11,728,328,854\\1040/1040 image pairs\\1040 size differences\end{tabular} \\
\hline
\end{tabular}
\par\smallskip
\begin{minipage}{0.98\textwidth}\footnotesize \textsuperscript{*}SheetShrink uses separately regenerated controls whose uncompressed members match the native candidates; full per-trial member inventories of the timed pipelines were not retained. Wanderungen uses a normalized baseline. Counts compare exact RGBA values in common rectangles with matched origins; overlapping captures can count pixels repeatedly. Image-size differences remain separate from common-area pixel counts. Different plans and incomplete scopes preclude a cross-tool preservation ranking. No capture regions or retained pairs provide no pixel-equality observation. Gray bars show the original; blue/orange show payload at or below/above the original. Input/output hashes, scope details and failures remain in the artifact.\end{minipage}
\end{table*}

\begin{table*}[p]
\centering
\footnotesize
\setlength{\tabcolsep}{2.5pt}
\renewcommand{\arraystretch}{1.08}
\caption{Opening/saving observations on two hosts. Cells show median total seconds above maximum memory across three attempts. Host blocks specify commit caps and memory metrics. The timeout is 900 s.}
\label{tab:case-study-resources}
\begin{tabular}{lrrrrrr}
\hline
Endpoint & Original & SheetShrink & NXPowerLite & XLStylesTool & \shortstack{Check\\Performance} & Inquire \\
\multicolumn{7}{l}{\textbf{Ryzen 5 4500U laptop: 4 GiB commit cap; library working set (MiB)}} \\
\hline
\multicolumn{7}{l}{\textbf{Destatis Migration}} \\
openpyxl & \shortstack[r]{\CasePairBars[2.2]{111.30}{1.10000}{1.10000}{casebarblue}\\\CasePairBars[2.2]{2700}{1.09997}{1.09997}{casebarblue}} & \shortstack[r]{\CasePairBars[2.2]{9.67}{1.10000}{0.09555}{casebarblue}\\\CasePairBars[2.2]{250}{1.09997}{0.10194}{casebarblue}} & \shortstack[r]{\CasePairBars[2.2]{83.58}{1.10000}{0.82602}{casebarblue}\\\CasePairBars[2.2]{1996}{1.09997}{0.81334}{casebarblue}} & \shortstack[r]{\CasePairBars[2.2]{111.13}{1.10000}{1.09829}{casebarblue}\\\CasePairBars[2.2]{2700}{1.09997}{1.10000}{casebarorange}} & \shortstack[r]{\CasePairBars[2.2]{9.11}{1.10000}{0.09002}{casebarblue}\\\CasePairBars[2.2]{236}{1.09997}{0.09603}{casebarblue}} & \shortstack[r]{\CasePairBars[2.2]{54.96}{1.10000}{0.54318}{casebarblue}\\\CasePairBars[2.2]{1530}{1.09997}{0.62347}{casebarblue}} \\
pandas to XLSX & \shortstack[r]{\colorbox{caseadverse}{\strut M}\\3900} & \shortstack[r]{35.83\\856} & \shortstack[r]{\colorbox{caseadverse}{\strut M}\\3900} & \shortstack[r]{\colorbox{caseadverse}{\strut M}\\3900} & \shortstack[r]{\colorbox{caseadverse}{\strut M}\\3900} & \shortstack[r]{\colorbox{caseadverse}{\strut M}\\3900} \\
pandas to CSV & \shortstack[r]{\colorbox{caseadverse}{\strut M}\\3900} & \shortstack[r]{11.54\\214} & \shortstack[r]{\colorbox{caseadverse}{\strut M}\\3900} & \shortstack[r]{\colorbox{caseadverse}{\strut M}\\3900} & \shortstack[r]{\colorbox{caseadverse}{\strut M}\\3900} & \shortstack[r]{\colorbox{caseadverse}{\strut M}\\3900} \\
Excel COM & \shortstack[r]{\CasePairBars[2.2]{20.34}{1.10000}{1.10000}{casebarblue}\\\CasePairBars[2.2]{429}{1.10000}{1.10000}{casebarblue}} & \shortstack[r]{\CasePairBars[2.2]{9.92}{1.10000}{0.53635}{casebarblue}\\\CasePairBars[2.2]{309}{1.10000}{0.79239}{casebarblue}} & \shortstack[r]{\CasePairBars[2.2]{17.94}{1.10000}{0.96972}{casebarblue}\\\CasePairBars[2.2]{409}{1.10000}{1.04743}{casebarblue}} & \shortstack[r]{\CasePairBars[2.2]{20.17}{1.10000}{1.09070}{casebarblue}\\\CasePairBars[2.2]{428}{1.10000}{1.09743}{casebarblue}} & \shortstack[r]{\CasePairBars[2.2]{5.34}{1.10000}{0.28867}{casebarblue}\\\CasePairBars[2.2]{259}{1.10000}{0.66287}{casebarblue}} & \shortstack[r]{\CasePairBars[2.2]{9.99}{1.10000}{0.53998}{casebarblue}\\\CasePairBars[2.2]{324}{1.10000}{0.83184}{casebarblue}} \\
R openxlsx & \shortstack[r]{\CasePairBars[2.2]{44.38}{1.08690}{1.08690}{casebarblue}\\\CasePairBars[2.2]{824}{1.09028}{1.09028}{casebarblue}} & \shortstack[r]{\CasePairBars[2.2]{13.70}{1.08690}{0.33547}{casebarblue}\\\CasePairBars[2.2]{317}{1.09028}{0.41938}{casebarblue}} & \shortstack[r]{\CasePairBars[2.2]{30.04}{1.08690}{0.73571}{casebarblue}\\\CasePairBars[2.2]{684}{1.09028}{0.90437}{casebarblue}} & \shortstack[r]{\CasePairBars[2.2]{44.92}{1.08690}{1.10000}{casebarorange}\\\CasePairBars[2.2]{831}{1.09028}{1.10000}{casebarorange}} & \shortstack[r]{\CasePairBars[2.2]{11.70}{1.08690}{0.28648}{casebarblue}\\\CasePairBars[2.2]{286}{1.09028}{0.37824}{casebarblue}} & \shortstack[r]{\CasePairBars[2.2]{32.22}{1.08690}{0.78900}{casebarblue}\\\CasePairBars[2.2]{564}{1.09028}{0.74618}{casebarblue}} \\
R readxl + writexl & \shortstack[r]{\colorbox{caseadverse}{\strut M}\\4070} & \shortstack[r]{178.17\\262} & \shortstack[r]{\colorbox{caseadverse}{\strut M}\\4070} & \shortstack[r]{\colorbox{caseadverse}{\strut M}\\4070} & \shortstack[r]{\colorbox{caseadverse}{\strut M}\\4072} & \shortstack[r]{\colorbox{caseadverse}{\strut M}\\4072} \\
\hline
\multicolumn{7}{l}{\textbf{Data.gov QT/RR}} \\
openpyxl & \shortstack[r]{\colorbox{caseadverse}{\strut M}\\3913} & \shortstack[r]{19.84\\318} & \shortstack[r]{12.26\\299} & \shortstack[r]{\colorbox{caseadverse}{\strut M}\\3913} & \shortstack[r]{12.07\\299} & \shortstack[r]{12.39\\303} \\
pandas to XLSX & \shortstack[r]{\CasePairBars[2.2]{249.13}{1.10000}{1.10000}{casebarblue}\\\CasePairBars[2.2]{327}{1.04035}{1.04035}{casebarblue}} & \shortstack[r]{\CasePairBars[2.2]{22.99}{1.10000}{0.10151}{casebarblue}\\\CasePairBars[2.2]{346}{1.04035}{1.10000}{casebarorange}} & \shortstack[r]{\CasePairBars[2.2]{15.79}{1.10000}{0.06972}{casebarblue}\\\CasePairBars[2.2]{327}{1.04035}{1.03958}{casebarblue}} & \shortstack[r]{\CasePairBars[2.2]{246.68}{1.10000}{1.08915}{casebarblue}\\\CasePairBars[2.2]{327}{1.04035}{1.03929}{casebarblue}} & \shortstack[r]{\CasePairBars[2.2]{15.55}{1.10000}{0.06864}{casebarblue}\\\CasePairBars[2.2]{327}{1.04035}{1.03987}{casebarblue}} & \shortstack[r]{\CasePairBars[2.2]{16.09}{1.10000}{0.07103}{casebarblue}\\\CasePairBars[2.2]{327}{1.04035}{1.03944}{casebarblue}} \\
Excel COM & \shortstack[r]{\CasePairBars[2.2]{39.79}{0.70004}{0.70004}{casebarblue}\\\CasePairBars[2.2]{929}{1.09853}{1.09853}{casebarblue}} & \shortstack[r]{\CasePairBars[2.2]{11.84}{0.70004}{0.20838}{casebarblue}\\\CasePairBars[2.2]{272}{1.09853}{0.32148}{casebarblue}} & \shortstack[r]{\CasePairBars[2.2]{4.67}{0.70004}{0.08223}{casebarblue}\\\CasePairBars[2.2]{253}{1.09853}{0.29933}{casebarblue}} & \shortstack[r]{\CasePairBars[2.2]{39.02}{0.70004}{0.68650}{casebarblue}\\\CasePairBars[2.2]{931}{1.09853}{1.10000}{casebarorange}} & \shortstack[r]{\CasePairBars[2.2]{62.52}{0.70004}{1.10000}{casebarorange}\\\CasePairBars[2.2]{246}{1.09853}{0.29069}{casebarblue}} & \shortstack[r]{\CasePairBars[2.2]{4.77}{0.70004}{0.08385}{casebarblue}\\\CasePairBars[2.2]{247}{1.09853}{0.29152}{casebarblue}} \\
R openxlsx & \shortstack[r]{\CasePairBars[2.2]{140.86}{1.10000}{1.10000}{casebarblue}\\\CasePairBars[2.2]{2453}{1.09990}{1.09990}{casebarblue}} & \shortstack[r]{\CasePairBars[2.2]{11.28}{1.10000}{0.08812}{casebarblue}\\\CasePairBars[2.2]{242}{1.09990}{0.10870}{casebarblue}} & \shortstack[r]{\CasePairBars[2.2]{4.26}{1.10000}{0.03324}{casebarblue}\\\CasePairBars[2.2]{224}{1.09990}{0.10044}{casebarblue}} & \shortstack[r]{\CasePairBars[2.2]{140.40}{1.10000}{1.09640}{casebarblue}\\\CasePairBars[2.2]{2454}{1.09990}{1.10000}{casebarorange}} & \shortstack[r]{\CasePairBars[2.2]{3.55}{1.10000}{0.02770}{casebarblue}\\\CasePairBars[2.2]{210}{1.09990}{0.09434}{casebarblue}} & \shortstack[r]{\CasePairBars[2.2]{4.32}{1.10000}{0.03375}{casebarblue}\\\CasePairBars[2.2]{225}{1.09990}{0.10065}{casebarblue}} \\
R readxl + writexl & \shortstack[r]{\CasePairBars[2.2]{11.90}{1.10000}{1.10000}{casebarblue}\\\CasePairBars[2.2]{1341}{1.10000}{1.10000}{casebarblue}} & \shortstack[r]{\CasePairBars[2.2]{10.06}{1.10000}{0.93020}{casebarblue}\\\CasePairBars[2.2]{258}{1.10000}{0.21204}{casebarblue}} & \shortstack[r]{\CasePairBars[2.2]{2.94}{1.10000}{0.27160}{casebarblue}\\\CasePairBars[2.2]{240}{1.10000}{0.19675}{casebarblue}} & \shortstack[r]{\CasePairBars[2.2]{11.89}{1.10000}{1.09963}{casebarblue}\\\CasePairBars[2.2]{1340}{1.10000}{1.09966}{casebarblue}} & \shortstack[r]{\CasePairBars[2.2]{2.88}{1.10000}{0.26619}{casebarblue}\\\CasePairBars[2.2]{248}{1.10000}{0.20355}{casebarblue}} & \shortstack[r]{\CasePairBars[2.2]{2.96}{1.10000}{0.27394}{casebarblue}\\\CasePairBars[2.2]{239}{1.10000}{0.19643}{casebarblue}} \\
\hline
\multicolumn{7}{l}{\textbf{Ryzen 7 7800X3D: 24 GiB commit cap; process-tree commit (GiB)}} \\
\multicolumn{7}{l}{\textbf{Destatis Migration}} \\
pandas to XLSX & \shortstack[r]{\colorbox{caseadverse}{\strut M}\\24.00} & \shortstack[r]{24.729\\1.27} & -- & -- & -- & -- \\
pandas to CSV & \shortstack[r]{\colorbox{caseadverse}{\strut M}\\24.00} & \shortstack[r]{8.291\\0.65} & -- & -- & -- & -- \\
R readxl + writexl & \shortstack[r]{392.149\\8.70} & \shortstack[r]{131.953\\0.51} & -- & -- & -- & -- \\
\hline
\multicolumn{7}{l}{\textbf{Data.gov QT/RR}} \\
openpyxl & \shortstack[r]{279.638\\12.56} & \shortstack[r]{11.519\\3.66} & -- & -- & -- & -- \\
\hline
\end{tabular}
\par\smallskip
\begin{minipage}{0.98\textwidth}\footnotesize SheetShrink includes fresh cleanup through library exit. Other columns measure processing of prepared outputs, with their cleanup excluded. Completed cells have three successful attempts. M means all three attempts stopped at the memory limit, with reached memory retained below. Bars compare original (gray) and measured value (blue at or below, orange above) on a common linear scale per endpoint and metric, requiring completed originals. RSS covers library phases, with nominal 20 ms sampling. Pipeline working sets, OS job-commit peaks, cleanup/open/save phases, stopped elapsed times, CPU, ranges and all 288 receipts remain in the artifact. The larger-host block reports aggregate process-tree commit in GiB, including cleanup. Dashes indicate unmeasured conditions.\end{minipage}
\end{table*}

